\scanpage{549}

and we assume that $\mathcal{G}^+\ne\mathcal{G}$, hence
\[
\tag{3}
\begin{gathered}
\mathcal{G}\simeq\{1,\tau'\}\mathbin{.}\mathcal{G}^+\\[-2pt]
(\tfrac12\ \text{direct})
\end{gathered}
\]
and we have a canonical exact sequence
\[
\tag{4}
\diagram{\begin{tikzcd}[column sep=large]
1\arrow[r]&\mathcal{G}^+\arrow[r]&\mathcal{G}\arrow[r,"\varepsilon_{\mathcal{G}}"]&\{\pm1\}\arrow[r]&1
\end{tikzcd}}
\]
We also denote by $\nu_{\mathcal{G}}$ the composite
\[
\tag{5}
\diagram{\begin{tikzcd}[column sep=large,row sep=small]
\mathcal{G}\arrow[r,"\varepsilon_{\mathcal{G}}"]\arrow[rr,bend left=25,"\nu_{\mathcal{G}}"]&\{\pm1\}&\mathbb{Z}/2\mathbb{Z}\arrow[l,"\sim"']\\
&(-1)^\nu&\nu\arrow[l,mapsto]
\end{tikzcd}}\ .
\]

In general, we \unsure{will denote} (by abuse of language) by the same letter elements of $\mathrm{Gl}(2,\mathbb{Z})$ --- such as $\varepsilon_0,\varepsilon_1,\tau=\tau_\infty,\tau'=\tau'_\infty$, $\sigma=\sigma_\infty,\rho=\rho_\infty$ etc --- and their images in $\mathcal{G}$.

\margin{The groups $Z_{\rho,\sigma}$, $\gamma_{\rho,\sigma}$}

$2^\circ)$ Let
\[
\tag{6}
\begin{gathered}
\deleted{B}Z_{\rho,\sigma}\ \subset\ \operatorname{Aut}(\mathcal{G}^+)\times\widehat{\mathbb{Z}}^*\times\{\pm1\}\times\{\pm1\}\\[-2pt]
\Vert\ \mathrm{dfn}\\[-2pt]
\left\{(u,\mu,\varepsilon,\varepsilon')\ \middle|\
\begin{aligned}
u(\varepsilon_0)&=\varepsilon_0^\mu\\
u(\rho)&\ \text{conjugate in }\mathcal{G}^+\text{ to }\rho^{\varepsilon'}\\
u(\sigma)&\ \rule{4.1cm}{0.3pt}\ \text{to }\sigma^\varepsilon
\end{aligned}\right\}
\end{gathered}
\]
This is a subgroup of the product group.

Let us put
\[
\tag{7}
\left\{\begin{aligned}
L_\rho&=\text{subgroup of }\mathcal{G}\text{ generated by }\rho\\
L_\sigma&=\rule{7cm}{0.3pt}\ \sigma
\end{aligned}\right.
\]
so that we have \unsure{epimorphisms}
\[
\tag{8}
\left\{\diagram{\begin{tikzcd}[column sep=large,row sep=small]
\mathbb{Z}/6\mathbb{Z}\arrow[r,"\text{epimorphism}"]&L_\rho&n\arrow[r,mapsto]&\rho^n\\
\mathbb{Z}/4\mathbb{Z}\arrow[r]&L_\sigma&n\arrow[r,mapsto]&\unsure{\sigma^n}
\end{tikzcd}}\right.\ ,
\]
\mspage{548}

{\small
\begin{center}
548\\
\texttt{https://grothendieck.umontpellier.fr}\\
Université de Montpellier, UMR CNRS 5149, Case courrier 051, Place Eugène Bataillon\\
34095 Montpellier cedex 5 - France
\end{center}
}

\scanpage{550}
\mspage{578}

and \deleted{some} \inserted{a commutative\ diagram of }homomorphisms
\[
\tag{9}
\diagram{\begin{tikzcd}[column sep=large,row sep=small]
&L_\rho\arrow[dr,hook]&\\
\{\pm1\}\arrow[ur]\arrow[dr]&&\mathcal{G}^+\\
&L_\sigma\arrow[ur,hook]&
\end{tikzcd}}
\]

--- If $\mathcal{G}$ is ``big enough'', we will moreover have that $\{\pm1\}\hookrightarrow\mathcal{G}^+$ i.e. ``$-1$'' $\ne1$ in $\mathcal{G}^+$, and even
\[
L_\rho\cap L_\sigma=\{\pm1\}\ .
\]

The conditions (6) therefore imply that
\[
\tag{10}
\left\{\begin{aligned}
u(L_\rho)&\ \inserted{\mathcal{G}^+\text{-}}\text{conjugate to }L_\rho\\
u(L_\sigma)&\ \mathcal{G}^+\text{-conjugate to }L_\sigma
\end{aligned}\right.
\]

If $\mathcal{G}^+$ is ``big enough'' so that $\rho$ and $\rho^{-1}$, $\sigma$ and $\sigma^{-1}$ are not conjugate in it, then an element $(u,\mu,\varepsilon,\varepsilon')$ of $Z_{\rho,\sigma}$ is known from the sole knowledge of $(u,\mu)$. If on the other hand the homomorphism
\[
\tag{11}
\diagram{\begin{tikzcd}[column sep=large,row sep=small]
\widehat{\mathbb{Z}}\arrow[r]&\mathcal{G}^+\\
n\arrow[r,mapsto]&\varepsilon_0^n
\end{tikzcd}}
\]
is injective, then $\mu$ is known when we know $u$. Thus if $\mathcal{G}^+$ is big enough to satisfy the three conditions considered (for example so that we have a ``modular homomorphism''
\[
\tag{12}
\diagram{\begin{tikzcd}[column sep=large]
\mathcal{G}\arrow[r]&\mathrm{Gl}(2,\widehat{\mathbb{Z}})
\end{tikzcd}}
\]
(sending $\rho,\sigma,\tau'$ to the elements with the same names), then an element of $Z_{\rho,\sigma}$ is known from the sole
\mspage{549}

{\small
\begin{center}
549\\
\texttt{https://grothendieck.umontpellier.fr}\\
Université de Montpellier, UMR CNRS 5149, Case courrier 051, Place Eugène Bataillon\\
34095 Montpellier cedex 5 - France
\end{center}
}

\scanpage{551}

knowledge of $u\in\operatorname{Aut}(\mathcal{G}^+)$.

We have a homomorphism
\[
\tag{13}
\diagram{\begin{tikzcd}[column sep=large,row sep=small]
L_0\arrow[r,"i_{L_0,Z}"]&\unsure{Z_{\rho,\sigma}}\\
\ell\arrow[r,mapsto]&(\operatorname{int}(\ell),1,1,1)
\end{tikzcd}}
\]
whose image is \unsure{normal}, the kernel is none other than $L_0\cap\operatorname{Cent}(\mathcal{G}^+)$. Let us note that we have
\[
\tag{14}
L_0\cap\operatorname{Cent}(\mathcal{G}^+)\subset L_0\cap\operatorname{Cent}_{\mathcal{G}^+}(\rho)\subset L_0\cap L_1
\]
for if $\ell\in L_0$ and $\rho$ commute, \inserted{i.e. \uline{if $\ell$}$=\rho(\ell)$,} since \deleted{$\ell=\rho(\ell)$} $\rho(L_0)=L_1$, we will have $\ell\in L_1$. Thus if $L_0\cap L_1=\{1\}$ (which is the case if (12) exists), \deleted{then} and even if there exists only
\[
\tag{12'}
\diagram{\begin{tikzcd}[column sep=large]
\mathcal{G}\arrow[r]&\mathrm{Gl}(2,\widehat{\mathbb{Z}})'
\end{tikzcd}}\qquad\text{modular},
\]
then we will have
\[
\tag{15}
L_0\cap\operatorname{Cent}(\mathcal{G}^+)=\{1\},
\]
and (13) will be injective. We are going to assume this, to fix ideas.

We put
\[
\tag{16}
\gamma_{\rho,\sigma}=Z_{\rho,\sigma}/\operatorname{Im}L_0,
\]
so that we will have a canonical exact sequence
\[
\tag{17}
\diagram{\begin{tikzcd}[column sep=large]
1\arrow[r]&L_0\arrow[r,"i_{L_0,Z}"]&Z_{\rho,\sigma}\arrow[r,"\delta_Z"]&\gamma_{\rho,\sigma}\arrow[r]&1
\end{tikzcd}}\ .
\]

The canonical projection homomorphisms $(u,\mu,\varepsilon,\varepsilon')\mapsto\mu$, $\varepsilon,\varepsilon'$ factor through $\gamma_{\rho,\sigma}$ and give canonical homomorphisms
\mspage{550}

{\small
\begin{center}
550\\
\texttt{https://grothendieck.umontpellier.fr}\\
Université de Montpellier, UMR CNRS 5149, Case courrier 051, Place Eugène Bataillon\\
34095 Montpellier cedex 5 - France
\end{center}
}

\scanpage{552}
\mspage{579}

\[
\tag{18}
\left\{\diagram{\begin{tikzcd}[column sep=large,row sep=small]
\gamma_{\rho,\sigma}\arrow[r,"\chi_\gamma"]&\widehat{\mathbb{Z}}^*\\
\gamma_{\rho,\sigma}\arrow[r,"\varepsilon_\gamma"]&\{\pm1\}\\
\gamma_{\rho,\sigma}\arrow[r,"\varepsilon'_\gamma"]&\{\pm1\}
\end{tikzcd}}\right.
\quad\text{i.e.}\quad
\diagram{\begin{tikzcd}[column sep=huge]
\gamma_{\rho,\sigma}\arrow[r,"(\chi_\gamma,\varepsilon_\gamma,\varepsilon'_\gamma)"]&\widehat{\mathbb{Z}}^*\times\{\pm1\}\times\{\pm1\}
\end{tikzcd}}
\]
the composites of these with $Z_{\rho,\sigma}\to\gamma_{\rho,\sigma}$ are denoted
\[
\tag{19}
\chi_Z=\chi_\gamma\circ\delta_Z,\quad
\varepsilon_Z=\varepsilon_\gamma\circ\delta_Z,\quad
\varepsilon'_Z=\varepsilon'_\gamma\circ\delta_Z,
\]
so that
\[
\tag{20}
\chi_Z(u,\mu,\varepsilon,\varepsilon')=\mu,\quad
\varepsilon_Z(u,\mu,\varepsilon,\varepsilon')=\varepsilon,\quad
\varepsilon'_Z(u,\mu,\varepsilon,\varepsilon')=\varepsilon'
\]

We would feel like denoting by $\gamma_{\rho,\sigma}^{\unsure{0|\circ}}$ the kernel of $(\varepsilon_\gamma,\varepsilon'_\gamma)$, by $\gamma_{\rho,\sigma}^{+}$ that of $\chi_\gamma$, by $\gamma_{\rho,\sigma}^{\unsure{0|\circ}+}$ their intersection, and introducing under the notations $Z_{\rho,\sigma}^{\unsure{0|\circ}}$, $Z_{\rho,\sigma}^{+}$, $Z_{\rho,\sigma}^{\unsure{0|\circ}+}$ their inverse images in $Z_{\rho,\sigma}$, i.e. the kernels \unsure{resp.$^{\mathrm t}$|resp.} of $(\varepsilon_Z,\varepsilon'_Z)$, of $\chi_Z$ and of $(\chi_Z,\varepsilon_Z,\varepsilon'_Z)$.

\margin{The groups $S_0G_{\rho,\sigma}$, $G_{\rho,\sigma}$}

$3^\circ)$ \deleted{The group $Z_{\rho,\sigma}$ \unsure{operates},} \inserted{We have a tautological homomorphism}
\[
\tag{21}
\diagram{\begin{tikzcd}[column sep=large,row sep=small]
Z_{\rho,\sigma}\arrow[r]&\operatorname{Aut}(\mathcal{G}^+)\\
(u,\mu,\varepsilon,\varepsilon')\arrow[r,mapsto]&u
\end{tikzcd}}\ ,
\]
thus $Z_{\rho,\sigma}$ \unsure{acts|operates} on $\mathcal{G}^+$, and we can therefore consider the $\tfrac12$ direct product
\[
\tag{22}
\begin{gathered}
S_0G_{\rho,\sigma}=Z_{\rho,\sigma}\mathbin{.}\mathcal{G}^+\ .\\[-2pt]
(\tfrac12\text{-direct})
\end{gathered}
\]
\mspage{551}

{\small
\begin{center}
551\\
\texttt{https://grothendieck.umontpellier.fr}\\
Université de Montpellier, UMR CNRS 5149, Case courrier 051, Place Eugène Bataillon\\
34095 Montpellier cedex 5 - France
\end{center}
}

\scanpage{553}

We have a homomorphism
\[
\tag{23}
\diagram{\begin{tikzcd}[column sep=large,row sep=small]
L_0\arrow[r,"i_{L_0,SG}"]&S_0G_{\rho,\sigma}=Z_{\rho,\sigma}\mathbin{.}\mathcal{G}^+\\
\ell\arrow[r,mapsto]&i_{L_0,Z}(\ell)\mathbin{.}\ell^{-1}
\end{tikzcd}}
\]
whose image is again normal, \deleted{\illegible{1}} \inserted{We} \deleted{\unsure{is stable}} \inserted{\unsure{put}}
\[
\tag{24}
G_{\rho,\sigma}=\underbrace{S_0G_{\rho,\sigma}}_{Z_{\rho,\sigma}\mathbin{.}\mathcal{G}^+}/\operatorname{Im}L_0\ \deleted{\unsure{Z_\rho|\mathcal{G}}},
\]
so that we will have an exact sequence
\[
\tag{25}
\diagram{\begin{tikzcd}[column sep=large]
1\arrow[r]&L_0\arrow[r,no head,"i_{L_0,SG}"]&S_0G_{\rho,\sigma}\arrow[r]&G_{\rho,\sigma}\arrow[r,no head]&1
\end{tikzcd}}\ ,
\]
and \unsure{a} diagram \inserted{\unsure{commutative}} of homomorphisms
\[
\tag{26}
\diagram{\begin{tikzcd}[column sep=large,row sep=small]
&Z_{\rho,\sigma}\arrow[dr,hook,"i_Z"]&\\
L_0\arrow[ur,hook,"i_{L_0,Z}"]\arrow[dr,hook,"i_{L_0,\mathcal{G}}"']&&G_{\rho,\sigma}\\
&\mathcal{G}^+\arrow[ur,hook,"i_{\mathcal{G}^+}"']&
\end{tikzcd}}
\]
(a priori $\operatorname{Ker}(\mathcal{G}^+\to G_{\rho,\sigma})\simeq\operatorname{Ker}(L_0\to Z_{\rho,\sigma})\simeq L_0\cap\operatorname{Cent}(\mathcal{G}^+)$, but by hypothesis (15) it is reduced to $\{1\}$), fitting into a commutative diagram of exact sequences
\[
\tag{27}
\diagram{\begin{tikzcd}[column sep=large,row sep=large]
1\arrow[r]&L_0\arrow[r,"i_{L_0,Z}"]\arrow[d,hook,"i_{L_0,\mathcal{G}}"']&Z_{\rho,\sigma}\arrow[r,"\delta_Z"]\arrow[d,hook]&\gamma_{\rho,\sigma}\arrow[r]\arrow[d,equal]&1\\
1\arrow[r]&\mathcal{G}^+\arrow[r,"i_{\mathcal{G}^+}"]&G_{\rho,\sigma}\arrow[r,"\delta_G"]&\gamma_{\rho,\sigma}\arrow[r,no head]&1
\end{tikzcd}}\ .
\]
We will put
\[
\tag{28}
\chi_G=\chi_\gamma\circ\delta_G,\quad
\varepsilon_G=\varepsilon_\gamma\circ\delta_G,\quad
\varepsilon'_\gamma\circ\delta_G
\]
\mspage{552}

{\small
\begin{center}
552\\
\texttt{https://grothendieck.umontpellier.fr}\\
Université de Montpellier, UMR CNRS 5149, Case courrier 051, Place Eugène Bataillon\\
34095 Montpellier cedex 5 - France
\end{center}
}

\scanpage{554}
\mspage{580}

hence a homomorphism
\[
\tag{29}
\diagram{\begin{tikzcd}[column sep=large]
G_{\rho,\sigma} \arrow[r,"{(\chi_G,\varepsilon_G,\varepsilon'_G)}"] &
\widehat{\mathbb Z}^{*}\times\{\pm1\}\times\{\pm1\}
\end{tikzcd}},
\]
leading us to introduce $G_{\rho,\sigma}^{+}$, $G_{\rho,\sigma}^{\unsure{0}}$,
$G_{\rho,\sigma}^{+0}$ as usual.

The injection $i_{\mathcal G^+}:\mathcal G^+\hookrightarrow G_{\rho,\sigma}$
can be extended to an embedding \deleted{\illegible{1}}
(denoted $i_{\mathcal G}$)
\[
\tag{30}
\diagram{\begin{tikzcd}
i_{\mathcal G}:\mathcal G \arrow[r,hook] & G_{\rho,\sigma}
\end{tikzcd}},
\]
which I define by its value on \deleted{$\unsure{\tau'}$}
\inserted{$\tau=\sigma\tau'$}, by
\[
\tag{31}
i_{\mathcal G}(\tau\deleted{'})=
\deleted{\illegible{1}}\,i_Z(\tau_Z)
\]
\margin{If $\tau_\gamma$ denotes the image of $\tau=\tau_G\in G_{\rho,\sigma}$
in $\gamma_{\rho,\sigma}=G_{\rho,\sigma}/\mathcal G^+$,
$\mathcal G$ is identified with the inverse image in $G_{\rho,\sigma}$ of the subgroup
$\{1,\tau_\gamma\}$ of $\gamma_{\rho,\sigma}$.}
\[
\tag{32}
\left\{\begin{aligned}
&\tau_Z\in Z_{\rho,\sigma}\quad\text{defined as}\\
&\tau_Z=\left(
\underbrace{\operatorname{int}(\tau)|\mathcal G^+}_{\in\operatorname{Aut}(\mathcal G^+)},
\underset{\substack{\rotatebox{90}{$\in$}\\\widehat{\mathbb Z}^{*}}}{-1},
\underset{\substack{\rotatebox{90}{$\in$}\\\{\pm1\}}}{-1},
\underset{\substack{\rotatebox{90}{$\in$}\\\{\pm1\}}}{-1}\right).
\end{aligned}\right.
\]

From now on we identify $\mathcal G$ with a subgroup of $G_{\rho,\sigma}$,
and we denote by the same letter elements of $\mathcal G$
(such as $\tau,\tau',\rho,\sigma,\varepsilon_0,\varepsilon_1$, etc)
and their images in $G_{\rho,\sigma}$.
\margin{Note that $\tau'_G\notin Z_{\rho,\sigma}$ on the other hand we have
\deleted{$\unsure{\tau'_G}$} $\tau_G\in Z_{\rho,\sigma}$
\deleted{\illegible{1}}.}
We will moreover have
\[
\tag{32bis}
\chi_G|\mathcal G=\varepsilon_G|\mathcal G=\varepsilon'_G|\mathcal G
=\varepsilon_{\mathcal G}:\left\{\begin{array}{l}
\text{trivial on }\mathcal G^+\\
-1\text{ on }\mathcal G^-=\mathcal G\setminus\mathcal G^+.
\end{array}\right.
\]

\margin{The groups $Z_\rho,Z_\sigma,N_\rho,N_\sigma,SN_\rho,SN_\sigma$\unsure{\dots}}
$4^\circ)$ We \deleted{define} note that the element $\tau'$ of
\deleted{$\unsure{G_{\rho,\sigma}}$} $\mathcal G\subset G_{\rho,\sigma}$
normalizes $L_\rho$ and $L_\sigma$, therefore also normalizes
\[
\tag{33}
\left\{\begin{aligned}
Z_\rho&\overset{\mathrm{dfn}}=\operatorname{Cent}_{G_{\rho,\sigma}}(\rho)
=\operatorname{Cent}_{G_{\rho,\sigma}}(L_\rho),\\
Z_\sigma&\overset{\mathrm{dfn}}=\operatorname{Cent}_{G_{\rho,\sigma}}(\sigma)
=\operatorname{Cent}_{G_{\rho,\sigma}}(L_\sigma)
\end{aligned}\right.
\]
\mspage{553}

{\small
\begin{center}
553\\
\texttt{https://grothendieck.umontpellier.fr}\\
Université de Montpellier, UMR CNRS 5149, Case courrier 051, Place Eugène Bataillon\\
34095 Montpellier cedex 5 - France
\end{center}
}